\subsection{Hypergraph XOR systems}

The parity machinery of this chapter fixes pair contexts, and the
all-versus-nothing arguments of GHZ and the magic square live on
three-element contexts. The constraint shape lifts: an XOR system assigns
each context, of any arity, a required XOR target, and a global section is
a solution of the induced linear system over the two-element field. The
refutation certificate is that system's dual witness.

\begin{theorem}[The all-versus-nothing certificate]\label{thm:hypergraph-avn-certificate}\lean{Obsmetrologia.QuantumContextuality.not_hasGlobalSection_of_avn}\leanok\uses{thm:possibilistic-model}
A family of contexts of an XOR system, repetition allowed, in which every
measurement has even incidence while the targets XOR to true refutes every
global section: summing the constraints over the family yields \(0 = 1\)
over \(\mathbb{Z}/2\). The even-incidence cancellation
(\texttt{xorOver\_eq\_false\_of\_even\_incidence}) and the certificate
theorem depend on no axioms.
\end{theorem}

\begin{theorem}[Pair parity models are XOR systems]\label{thm:parity-pair-weld}\lean{Obsmetrologia.QuantumContextuality.parityModel_toXor_hasGlobalSection_iff}\leanok\uses{thm:hypergraph-avn-certificate, thm:pr-cohomological-contextuality}
For a parity model whose contexts are all pairs, the model and its XOR
system have the same global sections, and a 4-cycle with target sum true
is an all-versus-nothing certificate on the flattened cycle
(\texttt{not\_hasGlobalSection\_of\_pair\_cycle}) --- the chapter's cycle
criterion recovered at the hypergraph level. The PR box's certificate is
its four contexts (\texttt{prXor\_no\_global}).
\end{theorem}

\begin{theorem}[GHZ and the magic square are strongly contextual as XOR systems]\label{thm:ghz-magic-square-avn}\lean{Obsmetrologia.QuantumContextuality.ghzXor_strongly_contextual}\leanok\uses{thm:hypergraph-avn-certificate, thm:magic-square-strong-contextual}
The GHZ system --- three parties carried on
\(\textup{Option Bool}\times\textup{Bool}\), contexts XXX, XYY, YXY, YYX ---
and the Mermin--Peres magic-square system on
\(\textup{Option Bool}\times\textup{Option Bool}\)
(\texttt{msXor\_strongly\_contextual}) are strongly contextual, each by one
certificate: all of its contexts. GHZ's contexts are triples
(\texttt{ghz\_not\_pairs}), outside the pair hypothesis of the parity
machinery. Both developments depend on no axioms. The identification of
the \(\textup{Fin } 9\) magic-square support model of
\cref{thm:magic-square-strong-contextual} with the XOR-system formulation
is an open weld.
\end{theorem}

\begin{theorem}[The Fredholm alternative over \(\mathbb{Z}/2\)]\label{thm:xor-fredholm-completeness}\notready\uses{thm:hypergraph-avn-certificate}
An XOR system with no all-versus-nothing certificate has a global section:
Gaussian elimination over the two-element field either solves the system
or produces a certificate from the eliminated rows. With
\cref{thm:hypergraph-avn-certificate} this makes strong contextuality of a
finite XOR system decidable by one rank computation, subsuming the pair,
GHZ, Mermin, and magic-square cases.
\end{theorem}

\begin{theorem}[The record mass of a hypergraph system]\label{thm:hypergraph-record-mass}\notready\uses{thm:record-merge-ground-state, thm:ghz-magic-square-avn}
The ground-state identity of \cref{thm:record-merge-ground-state} extends
past pair contexts: over an XOR system, the minimum total contradiction
mass of the counting merge over realizing families equals the minimum
frustrated-context count over the certificate coset. On GHZ the three-way
contexts price frustration differently than any pair scenario --- whether
the mass separates GHZ from the PR box is part of the claim.
\end{theorem}

\begin{theorem}[No proper sub-family of GHZ's four contexts is a certificate]\label{thm:ghz-proper-submask}\lean{Obsmetrologia.QuantumContextuality.ghz_proper_submask_not_avn, Obsmetrologia.QuantumContextuality.ghz_proper_submask_has_section}\leanok\uses{thm:ghz-magic-square-avn, thm:hypergraph-avn-certificate}
The finite combinatorial form of the Borromean reading, over all declared requiredteen sub-masks of GHZ's four
contexts (\texttt{ghzMask}). \texttt{ghz\_proper\_submask\_not\_avn} is proved from the design's
incidence structure rather than by enumerating the sixteen masks: GHZ's four contexts form a
2-regular, pairwise-singleton-intersecting block design (four blocks of size three on six points,
each point in exactly two blocks, each pair of blocks meeting in exactly one point), so a
measurement's incidence count in a sub-mask is exactly two minus the number of its two containing
contexts absent from the mask --- odd iff exactly one is absent. Four fixed measurements
(\texttt{ghzW1}--\texttt{ghzW4}, one per pair of contexts they jointly occur in) have this parity
computed once each (\texttt{ghzWk\_parity\_eq}), and a four-bit disagreement lemma
(\texttt{ghz\_bits\_disagree\_of\_not\_all\_eq}: bits not all equal disagree at one of four fixed
pairs) shows every mask other than the all-false and all-true ones fails to be even everywhere at
one of the four witnesses (\texttt{ghz\_proper\_submask\_odd\_witness}); the all-false mask is
handled directly (empty family, target sum \texttt{false}), and the all-true mask is excluded by
hypothesis. The only closed enumeration remaining is \texttt{ghz\_bits\_agree\_forces\_boundary},
a sixteen-value numeral fact (which \(n<16\) have all four bits equal) --- an arithmetic check, not
a search over GHZ's combinatorics. \texttt{ghz\_proper\_submask\_has\_section}: every proper
sub-family has a genuine satisfying assignment on its surviving measurements, found by a bounded
search over the six measurements' \(2^6\) assignments (not yet replaced by a constructive
per-case witness). Together: GHZ's obstruction is invisible to any proper sub-family. All carry
\texttt{propext}; the incidence lemmas and the section theorem also \texttt{Quot.sound}.
\end{theorem}

\begin{theorem}[The Alexander--Whitney cup product on \texttt{Cochain M}]\label{thm:cochain-cup-product}\lean{Obsmetrologia.QuantumContextuality.cupAt, Obsmetrologia.QuantumContextuality.deltaX_cupAt_one, Obsmetrologia.QuantumContextuality.deltaX_cupAt_one_leibniz}\leanok\uses{thm:hypergraph-xor-cochain-complex}
\texttt{cupAt p $\alpha$ $\beta$}: the front-\((p{+}1)\)/back-from-\(p\) split-and-\texttt{AND}
product, the \(\mathbb{Z}/2\) instance of the classical Alexander--Whitney cup product
\((\alpha\cup\beta)(v_0,\dots,v_{p+q}) = \alpha(v_0,\dots,v_p)\cdot\beta(v_p,\dots,v_{p+q})\).
For \(p=q=1\) the cochain-level Leibniz rule is proved by direct \texttt{deltaX\_cons} unfolding
on an explicit 4-vertex list (\texttt{deltaX\_cupAt\_one}, axiom-free), then regrouped
(\texttt{deltaX\_cupAt\_one\_leibniz}, axiom-free) into
\(\alpha[i,j]\wedge(\beta[j,k]\mathrel{\mathtt{xor}}\beta[j,l])$ xor $(\alpha[i,k]\mathrel{\mathtt{xor}}\alpha[j,k])\wedge\beta[k,l]\) --- a
front-face term reading $\beta$'s local variation against a fixed $\alpha$ edge, xored with a
back-face term reading $\alpha$'s variation against a fixed $\beta$ edge, the $\mathbb{Z}/2$
shape of $\delta(\alpha\cup\beta)=\alpha\cup\delta\beta \;\oplus\; \delta\alpha\cup\beta$. A first
attempt at this identity, built from a cup formula for two \emph{0}-cochains
misapplied to pair-valued data, failed a direct symbolic check; the general \((p,q)=(1,1)\)
telescoping identity was re-derived and confirmed against three independent published statements
of the classical formula before being proved in Lean, and the $p=q=0$ case (genuinely pointwise,
$(\alpha\cup\beta)(v)=\alpha(v)\cdot\beta(v)$, not pair-valued) was the source of the earlier
error. General $p,q$ and the full \(\delta\circ\delta\)-compatible cup-product axiom system
remain open past this one instance.
\end{theorem}

\begin{theorem}[GHZ's obstruction is a triple product, not a pairwise one]\label{thm:ghz-massey-triple}\notready\uses{thm:ghz-proper-submask, thm:cochain-cup-product}
The cohomological reading of \cref{thm:ghz-proper-submask}: the target identification with a
Massey product $\langle\alpha,\beta,\gamma\rangle$ on \texttt{Cochain GhzM}, where the three
pairwise cup products $\alpha\cup\beta$, $\beta\cup\gamma$, $\alpha\cup\gamma$ vanish (read off
\cref{thm:ghz-proper-submask}'s sections) while the triple product does not (read off its
certificate). \cref{thm:cochain-cup-product} now supplies a cup product and its Leibniz law for
the \(p=q=1\) case GHZ's classes would need, but the Massey product itself --- the secondary
operation defined on cohomology classes whose pairwise cup products are exact, valued in the
quotient by indeterminacy --- is not yet constructed; this node still reduces to the finite
combinatorial shadow of \cref{thm:ghz-proper-submask} until that construction exists.
\end{theorem}

\begin{theorem}[A cochain complex for hypergraph XOR systems]\label{thm:hypergraph-xor-cochain-complex}\lean{Obsmetrologia.QuantumContextuality.deltaX_deltaX, Obsmetrologia.QuantumContextuality.deltaX_single, Obsmetrologia.QuantumContextuality.coboundary_cocycle}\leanok\uses{thm:hypergraph-avn-certificate}
The carrier for \cref{thm:ghz-massey-triple}: \texttt{Cochain M} (\(\mathtt{List}\,M \to
\mathtt{Bool}\)) grades by context arity exactly as \texttt{XorSystem.target} already does, with
a general coboundary \texttt{deltaX} in the same cone-contraction form as
\texttt{CausalSets.Causal.Cohomology.N.delta}, \(\mathtt{xor}\) standing in for integer
subtraction: removing the head measurement, recursing into the tail with the prepend-shift.
\texttt{deltaX\_deltaX} proves \(\delta\circ\delta=0\) in every degree by the same chain-length
induction as the causal \(H^n\) module, with a pointwise extensionality lemma
(\texttt{deltaX\_congr}) standing in for \texttt{funext} at the prepend-shift step --- both depend
on no axioms, an improvement on the integer-coefficient original's recorded
\texttt{propext}/\texttt{Quot.sound} debt. The degree-one weld
(\texttt{deltaX\_single}): \texttt{deltaX} applied to a degree-0 cochain's singleton-list embedding,
evaluated on a pair, is \(\delta^0\) on that pair (\cref{thm:parity-pair-weld}'s \texttt{delta0},
up to \texttt{delta0\_symm}) --- not a definitional identity, since \texttt{Coboundary} above is
indexed by ordered pairs \(M\to M\to\mathtt{Bool}\) while \texttt{Cochain M} grades by arity, but
a proved specialization; this one lemma carries \texttt{propext}, from the singleton-list
match-definition. \texttt{coboundary\_cocycle} bundles \(\operatorname{im}\delta\subseteq
\ker\delta\), so a graded \(H^n\) over \(\mathbb{Z}/2\) is well-defined at every context arity,
including the triple-context degree GHZ and the magic square occupy. Still open: reading the
all-versus-nothing certificate of \cref{thm:hypergraph-avn-certificate} as a statement about
\(H^1\)/\(H^2\) of this complex, and weighing \cref{thm:hypergraph-record-mass} as a graded
invariant --- this node supplies the complex only, not yet the identification.
\end{theorem}
